\subsection{DERIVATIVES OF REAL-VALUED FUNCTIONS}

The preceding definitions and propositions may be augmented in several respects when we deal with real-valued functions of a real variable.

In the first place, if \(f\) is such a function, defined on an interval \(I \subset \mathbf{R}\) , and continuous relative to \(I\) at a point \(x_0 \in I\) , it can happen that when \(x\) tends to \(x_0\) while remaining in \(I\) and \(\neq x_0\) , that \(\frac{f(x) - f(x_0)}{x - x_0}\) has a limit equal to \(+\infty\) or to \(-\infty\) ; one then says that \(f\) is differentiable at \(x_0\) and has derivative \(+\infty\) (resp. \(-\infty\) ) there; if the function \(f\) has a derivative \(f'(x)\) (finite or infinite) at every point \(x\) of \(I\) , then the function \(f'\) (with values in \(\overline{\mathbf{R}}\) ) is again called the derived function (or simply the derivative) of \(f\) . One generalizes the definitions of right and left derivative similarly.

Example. At the point x = 0 the function \(x^{1/3}\) (the inverse function of \(x^{3}\) , a homeomorphism of R onto itself) has a derivative, equal to \(+\infty\) ; at x = 0 the function \(|x|^{1/3}\) has right derivative \(+\infty\) and left derivative \(-\infty\) .

The formulae for the derivative of a sum, of a product of differentiable real functions, and for the inverse of a differentiable function (props. 1, 3 and 4), as well as for the derivative of a (real-valued) composition of functions (prop. 5) remain valid when the derivatives that occur are infinite, so long as all the expressions that occur in these formulae make sense (Gen.~Top., IV, p.~345--346). In fact, if in prop. 6 one supposes that f is strictly increasing (resp. strictly decreasing) and continuous on I, and if \(f'(x_{0}) = 0\) , then the inverse function g has a derivative equal to \(+\infty\) (resp. \(-\infty\) ) at the point \(y_{0} = f(x_{0})\) ; if \(f'(x_{0}) = +\infty\) (resp. \(-\infty\) ), then g has derivative 0. There are similar results for right and left derivatives, which we leave to the care of the reader.

Let C be the graph or representing curve of a finite real function f, the subset of the plane \(R^{2}\) formed by the points \((x, f(x))\) where x runs through the set where f is defined. If the function f has a finite right derivative at a point \(x_{0} \in I\) , then the half-line with origin at the point \(\mathbf{M}_{x_{0}} = (x_{0}, f(x_{0}))\) of C, and direction numbers \((1, f_{d}'(x_{0}))\) is called the right half-tangent to the curve C at the point \(M_{x_{0}}\) ; when \(f_{d}'(x_{0}) = +\infty\) (resp. \(f_{d}'(x_{0}) = -\infty\) ) we use the same terminology for the half-line with origin \(M_{x_{0}}\) and direction numbers \((0, 1)\) (resp. \((0, -1)\) ). In the same way one defines the left half-tangent at \(M_{x_{0}}\) when \(f_{g}'(x_{0})\) exists. With these definitions one can verify quickly that the angle which the right (resp. left) half-tangent makes with the abscissa is the limit of the angle made by this axis with the half-line originating at \(M_{x_{0}}\) and passing through the point \(\mathbf{M}_{x} = (x, f(x))\) of C, as x tends to \(x_{0}\) while remaining \textgreater{} \(x_{0}\) (resp. \textless{} \(x_{0}\) ).

One can also say that the right (resp. left) half-tangent is the limit of the half-line originating at \(M_{x_{0}}\) passing through \(M_{x}\) , on endowing the set of half-lines having the same origin with the quotient space topology \(C^{*}/R_{+}^{*}\) (Gen.~Top., VIII, p.~107).

If the two half-tangents exist at a point \(M_{x_{0}}\) of C, they are in opposite directions only when f has a derivative (finite or not) at the point \(x_{0}\) (assumed interior to I); they are identical only when \(f_{d}^{\prime}(x_{0})\) and \(f_{g}^{\prime}(x_{0})\) are infinite and of opposite sign. In these two cases we say that the line containing these two half-tangents is the tangent to C at the point \(M_{x_{0}}\) .

When the tangent at \(M_{x_{0}}\) exists it is the limit of the line passing through \(M_{x_{0}}\) and \(M_{x}\) as x tends to \(x_{0}\) remaining \(\neq x_{0}\) , the topology on the set of lines which pass through a given fixed point being that of the quotient space \(C^{*}/R^{*}\) (Gen.~Top., VIII, p.~114).

The concepts of tangent and half-tangent to a graph are particular cases of general concepts which will be defined in the part of this Series devoted to differentiable varieties.

DEFINITION 4. We say that a real function \(f\) , defined on a subset \(A\) of a topological space \(E\) , has a relative maximum (resp. strict relative maximum, relative minimum, strict relative minimum) at a point \(x_0 \in A\) , relative to \(A\) , if there is a neighbourhood \(V\) of \(x_0\) in \(E\) such that at every point \(x \in V \cap A\) distinct from \(x_0\) one has \(f(x) \leqslant f(x_0)\) (resp. \(f(x) < f(x_0)\) , \(f(x) \geqslant f(x_0)\) , \(f(x) > f(x_0)\) ).

It is clear that if f attains its least upper bound (resp. greatest lower bound) over A at a point of A, then it has a relative maximum (resp. relative minimum) relative to A at this point; the converse is of course incorrect.

Note that if \(B \subset A\) , and if f admits (for example) a relative maximum at a point \(x_{0} \in B\) relative to B, then f does not necessarily have a relative maximum relative to A at this point.

PROPOSITION 7. Let \(f\) be a finite real function, defined on an interval \(\mathbf{I} \subset \mathbf{R}\) . If \(f\) admits a relative maximum (resp. relative minimum) at a point \(x_0\) interior to \(\mathbf{I}\) , and has both right and left derivatives at this point, then one has \(f_d'(x_0) \leqslant 0\) and \(f_g'(x_0) \geqslant 0\) (resp. \(f_d'(x_0) \geqslant 0\) and \(f_g'(x_0) \leqslant 0\) ); in particular, if \(f\) is differentiable at the point \(x_0\) , then \(f'(x_0) = 0\) .

The proposition follows trivially from the definitions.

We can say further that if at a point \(x_{0}\) interior to I the function f is differentiable and admits a relative maximum or minimum, then the tangent to its graph is parallel to the abscissa. The converse is incorrect, as is shown by the example of the function \(x^{3}\) which has zero derivative at the point x = 0, but has neither relative maximum nor minimum at this point.

